\section*{Multivibratore PWM}

\noindent\colorbox{orange}{
  \begin{minipage}{1.0\linewidth}
    \subsection*{Richiesta (a)}
    Descrivere qualitativamente il funzionamento del circuito, eventualmente individuando nel
    circuito alcune strutture gi\'a note e mostrando che per $V_M = 0$ si comporti come un
    multivibratore astabile e con $\text{duty-cycle}=50\%$ (nell'ipotesi che i livelli di
    saturazione siano simmetrici).
  \end{minipage}
}
\subsection*{Risposta (a)}
\paragraph{\colorbox{blu}{Condensatore scarico}}
Durante questa fase l'ingresso invertente \'e messo a terra, l'op-amp si comporta come un \textit{trigger di Schmitt \underline{non}-invertente}.
La tensione rimane nel livello di saturazione.

\paragraph{\colorbox{blu}{Transizione}}
Appena il condensatore riceve abbastanza carica cos\'i che $V_+ = V_-$\footnotemark
\footnotetext{Si parla della differenza di potenziale ma siccome entrambi gli ingressi sono collegati a terra...}
allora l'op-amp inizia ad entrare in regime di operazione lineare e si comporta da \textit{amplificatore \underline{invertente}}.

\paragraph{\colorbox{blu}{Condensatore carico}}
Durante il regime di amplificazione invertente, la tensione si ``capovolge'' e l'op-amp inizia a scaricarsi fino a
ri-attraversare la fase di transizione per tornare nella regione di operazione non-lineare come un \textit{trigger di Schmitt \underline{non}-invertente}. \\

\noindent\colorbox{orange}{
  \begin{minipage}{1.0\linewidth}
    \subsection*{Richiesta (b)}
    Nell'ipotesi di idealit\'a dell'operazionale, determinare l'andamento temporale del segnale $V_-$
    in funzione di $V_M$ e per ciascuno stato di $V_{OUT}$ (suggerimento: per i processi di carica/scarica
    di $C_1$ si sostituisca la rete costituita dal partitore $V_{OUT}-R_3-R_4-V_M$ con il suo equivalente di
    Th\'evenin visto dai terminali del condensatore; si consideri inoltre che $R_2 \ll R_1 \iff |V_+|\ll V_{OH}$,
    per cui i termini esponenziali possono essere approssimati al primo ordine in serie di Taylor
    come …).
  \end{minipage}
}
\subsection*{Risposta (b)}
La tensione equivalente di Th\'evenin
\begin{equation*}
  \begin{gathered}
    \frac{V_{OUT} - V_-}{Z_{R_3}} = \frac{V_- - V_M}{Z_{R_4}} \\ \implies
    V_- = \frac{V_{OUT} + V_M}{2}
  \end{gathered}
\end{equation*}
l'impedenza equivalente di Th\'evenin
\begin{equation*}
  Z_{th} = Z_{R_3} || Z_{R_4} = 50R
\end{equation*}
la legge oraria della tensione $V_-$
\begin{equation*}
  \begin{gathered}
    Z_{C_1}i_{C_1} = \frac{V_{th} - V_-}{Z_{th}} \\ \implies
    C\dot{V}_- = \frac{V_{th} - V_-}{50 R} \\ \implies
    \tau\dot{V}_- + V_- = V_{th} \quad \tau \equiv 50RC \\ \implies
    V_-(t) = V_{th} + [V_-(t_0) - V_{th}]\exp\left(-\frac{t - t_0}{\tau}\right) \\
    V_-(t) = \frac{V_{OH} + V_M}{2} + \left[V_{thr-} - \frac{V_{OH} + V_M}{2}\right]\exp\left(-\frac{t - t_0}{\tau}\right) \\
    V_-(t) = \frac{V_{OL} + V_M}{2} + \left[V_{thr+} - \frac{V_{OL} + V_M}{2}\right]\exp\left(-\frac{t - t_0}{\tau}\right) \\
  \end{gathered}
\end{equation*}

\noindent\colorbox{orange}{
  \begin{minipage}{1.0\linewidth}
    \subsection*{Richiesta (c)}
    In funzione di $V_M$ (oppure di $V_M/V_{OH}$) determinare la durata degli stati di $V_{OUT}$ e ricavare
    frequenza e duty-cycle dello stesso segnale, confrontandoli con quanto misurato al punto 2a.
  \end{minipage}
}
\subsection*{Risposta (c)}
il semi-periodo
\begin{equation*}
  \begin{gathered}
    t_L = \tau\ln\left(\frac{V_{thr+} - V_{Th,-}}{V_{thr-} - V_{Th,-}}\right) \\
    t_H = \tau\ln\left(\frac{V_{Th,+} - V_{thr-}}{V_{Th,+} - V_{thr+}}\right)
  \end{gathered}
\end{equation*}
approssimazione per Taylor siccome $R_2 \ll R_1$
\begin{equation*}
  \begin{gathered}
    t_L =  \tau\ln\left(1 + \frac{V_{thr+} - V_{thr-}}{V_{thr-} - V_{Th,-}}\right) \approx \tau\frac{V_{thr+} - V_{thr-}}{V_{thr-} - V_{Th,-}} \\
    t_H =  \tau\ln\left(1 + \frac{V_{thr+} - V_{thr-}}{V_{Th,+} - V_{thr+}}\right) \approx \tau\frac{V_{thr+} - V_{thr-}}{V_{Th,+} - V_{thr+}}
  \end{gathered}
\end{equation*}
sotto ipotesi di simmetria nelle tensione di saturazione
\begin{equation*}
  \begin{gathered}
    t_L = \tau\frac{2V_{thr}}{V_{thr} - V_{Th,-}} = \tau\frac{2\beta}{\beta - V_{Th,-}/V_{OH}} = \tau\frac{4\beta}{1 - 2\beta - V_M/V_{OH}} \\
    t_H = \tau\frac{2V_{thr}}{V_{Th,+} - V_{thr}} = \tau\frac{4\beta}{1 + V_M/V_{OH} - 2\beta}
  \end{gathered}
\end{equation*}
il duty cycle
\begin{equation*}
  D = \frac{1}{1 + \frac{(1 - 2\beta) - V_M/V_{OH}}{(1 - 2\beta) + V_M/V_{OH}}}
\end{equation*}

\noindent\colorbox{orange}{
  \begin{minipage}{1.0\linewidth}
    \subsection*{Richiesta (d)}
    Discutere gli effetti delle caratteristiche non ideali dell'operazionale, ed in
    particolare del suo slew-rate, sul funzionamento del circuito, ad esempio evidenziando i suoi
    effetti sulle forme d'onda misurate.
  \end{minipage}
}
\subsection*{Risposta (d)}
Lo slew-rate del circuito fa si che la tensione in uscita dall'op-amp non commuti istantaneamente
dalla tensione di saturazione bassa $V_{OL}$ alla tensione di saturazione alta $V_{OH}$.
